\noindent\textbf{Application setting.}
\method{} targets services that update node classifiers as documents, products,
or other text-bearing entities arrive. Training retains compact examples after
raw snapshots are retired, while serving maintains the prediction graph.
Construction uses the current increment and historical memory. The
class-incremental evaluation uses the benchmark graph, while the chronological
evaluation uses the stage-available graph with recurring classes and arrived
endpoints.

\noindent\textbf{Conditions and scope.}
Each increment must provide features, text, construction labels, structural
relations, and retainable source texts. The method assumes compatible
structural and textual signals, consistent feature coordinates, encoder, and
class meanings across stages. Under these conditions, learned support and
class matching improve performance, aggregation through $Q$ with a bounded
residual is more effective than direct feature optimization on SportsFit, and
shared support reduces construction cost. The value of gradient matching
depends on the dataset and budget. The results concern supervised node
classification with stable labels. Severe drift, revised labels, indefinitely
long streams under a hard storage limit, and relation-dependent tasks require
new memory policies or access rules.

\noindent\textbf{Construction and reuse.}
Class-block storage and cache offloading are optional implementation choices.
The resulting memory can be reused across updates or student architectures.
The reported downstream reductions use precomputed text embeddings with the
encoder off the GPU, so they do not represent end-to-end cost. Construction
remains an upfront cost, fixed-rate storage grows with the stream, and
inference must use the current graph when historical nodes are unavailable.
